\subsection{\problemblock{*Integ} Data Block}\label{sec:block-integ}

Trajectories are computed by integrating the equations of motion with a fourth-order,
fixed-step Runge--Kutta numerical method. The integration step size, which controls
the accuracy of the calculation, is fixed within each trajectory segment. It must be
small enough to capture changes in aerodynamic and propulsive forces, but large
enough that computation time is not excessive.

The default step size is 0.10~s. It can be changed with an
\problemblock{*integ} block; for example,

\begin{lstlisting}[style=taosmanual]
*integ dt=0.25
\end{lstlisting}

requests a 0.25~s step. Typical values range from 0.01 to 0.5~s. The minimum
allowable integration step is $1\times10^{-6}$~s.

Output variables are saved at the beginning and end of each segment and at even
increments of the print interval within the segment. In addition to controlling
printout density, the print interval controls memory use because all trajectory output
variables are retained in memory at those times. Small print intervals therefore
increase memory requirements, and large intervals reduce them. The print interval
\statevar{dtprnt} defaults to 1.0~s. For example,

\begin{lstlisting}[style=taosmanual]
*integ dtprnt=0.5 dt=0.05
\end{lstlisting}

changes the print interval to 0.5~s and the integration step to 0.05~s.

\taossubhead{Guidance Time Constant}

Some guidance rules, including those in \cref{tab:flight-condition-guidance},
specify control variables indirectly through a desired flight condition. When the
actual condition differs from the desired condition, TAOS changes the free controls
so that the trajectory moves toward the requested state. The rate of that correction is
controlled by the guidance time constant \statevar{dtguid}.

A small \statevar{dtguid} produces large corrections in an attempt to reach the
desired condition quickly. A large value produces small corrections and gradual
convergence. Thus, the parameter behaves somewhat like a time constant: small
values may produce oscillations that do or do not damp, while large values tend to
produce exponential convergence but may correct too slowly.

\begin{lstlisting}[style=taosmanual]
*integ dtprnt=1.0 dt=0.20 dtguid=5.0
\end{lstlisting}

sets the time constant to 5~s. This is considered a large value and causes only small
flight-path corrections; they may be too small to achieve the desired conditions.
Experimentation is sometimes required. The default value of 1~s works well for
most guidance rules.
